\documentclass[10pt,twocolumn]{article}
\usepackage[margin=0.72in,columnsep=0.25in]{geometry}
\usepackage{amsmath,amssymb,amsthm,mathtools}
\usepackage{graphicx}
\usepackage{booktabs}
\usepackage{caption}
\usepackage{hyperref}
\usepackage{algorithm}
\usepackage[noend]{algorithmic}
\usepackage{xcolor}
\usepackage{times}
\usepackage{enumitem}
\usepackage{bm}
\usepackage{tcolorbox}

\setlist{nosep,leftmargin=*}
\newcommand{\placeholder}[1]{\textcolor{red!70!black}{\texttt{[#1]}}}
\newcommand{\R}{\mathbb{R}}
\newcommand{\E}{\mathbb{E}}
\newcommand{\Var}{\mathrm{Var}}
\newcommand{\diag}{\mathrm{diag}}
\newcommand{\sign}{\mathrm{sign}}
\newcommand{\MAD}{\mathrm{MAD}}
\newcommand{\clip}{\mathrm{clip}}
\newcommand{\bw}{\bm{w}}
\newcommand{\bmu}{\bm{\mu}}

\newtheorem{definition}{Definition}
\newtheorem{proposition}{Proposition}
\newtheorem{remark}{Remark}

\title{\textbf{PRISM: Predictive Returns and Inferred Structure\\for Modern Portfolios}\\[4pt]
\large Sharpe--Optimal Allocation under Forecast Uncertainty\\with Kantorovich Order Flow}

\author{Benjamin Coriat\\
ECSE 551 --- Machine Learning for Engineers, McGill University\\
\texttt{benjamin.coriat@mail.mcgill.ca}}
\date{April 2026}

\begin{document}
\maketitle

\begin{abstract}
Machine learning has dramatically improved financial return
prediction, yet translating forecasts into portfolios remains
fragile: optimisers amplify estimation noise, and na\"ive
diversification often wins.  We present PRISM, a framework that
closes this gap by treating forecast uncertainty as a first-class
object.  Alongside every return and correlation prediction, the
pipeline produces a structured five-dimensional error profile
that quantifies \emph{how much to trust it}.  These profiles feed
a per-asset confidence score, which drives a Normal--Inverse-Wishart
Bayesian update that shrinks unreliable forecasts toward the market
mean.  The resulting regularised inputs enter a Sharpe-maximising
convex programme whose solution is executed at minimum transaction
cost via Sinkhorn optimal transport.  We backtest the full pipeline
on 24 global equity indices and commodities over 2012--2025 at
weekly frequency, with all strategies vol-targeted to 15\% for
fair comparison.  We provide complete mathematical derivations,
pseudocode aligned with the implemented codebase, and placeholders
for empirical results pending execution.
\end{abstract}

%======================================================================
\section{Introduction}\label{sec:intro}
%======================================================================

Markowitz~\cite{markowitz1952} showed that the mean--variance
optimal portfolio satisfies $\bw^*\propto\Sigma^{-1}\bmu$,
where $\bmu$ is the expected return vector and $\Sigma$ the
covariance matrix.  The formula is elegant but dangerous in
practice: $\Sigma^{-1}$ amplifies noise in $\bmu$, so small
estimation errors produce wildly different---and often
pathological---portfolios~\cite{merton1980}.  Michaud~\cite{michaud1989}
called this ``estimation error maximisation.''
DeMiguel et~al.~\cite{demiguel2009} provided the sharpest
indictment: across 14 models and 7 datasets, na\"ive
equal-weight ($1/N$) diversification matched or beat every
optimised strategy after transaction costs, precisely because
optimisers consume estimation noise as if it were signal.

Modern machine learning has significantly improved forecast
accuracy.  Gu et~al.~\cite{gu2020empirical} showed that neural
networks and tree ensembles produce economically meaningful
out-of-sample $R^2$ for individual stock returns.  Yet the standard
practice remains to plug ML predictions directly into an optimiser
as if they were ground truth---ignoring the fact that even a
state-of-the-art model's $R^2$ rarely exceeds 1--2\% at the weekly
horizon.  The result is that better predictions do not
automatically yield better portfolios.

PRISM addresses this disconnect.  The central idea is to produce,
alongside every forecast, a structured \emph{error profile} that
quantifies how reliable it is.  Concretely, at each weekly
rebalancing date the pipeline delivers five tensors:
\begin{enumerate}
\item $\hat\bmu_{t+1}\in\R^n$: predicted return vector.
\item $E_\mu\in\R^{n\times 5}$: return error characterisation
  (bias, directional accuracy, dispersion, skewness, kurtosis).
\item $\hat C_{t+1}\in\R^{n\times n}$: predicted correlation matrix.
\item $E_C\in\R^{n\times n\times 4}$: correlation error tensor.
\item $V\in\R^{n\times 26}$: volatility history matrix.
\end{enumerate}
The allocation layer consumes these tensors through mathematically
principled channels: $E_\mu$ feeds a confidence score that drives
Bayesian shrinkage of $\hat\bmu$; $E_C$ and $V$ inform the
covariance estimate; and the resulting regularised inputs enter a
Sharpe-maximising convex programme.  The target portfolio is then
executed via Sinkhorn optimal transport to minimise transaction
costs.

The paper is organised to mirror the pipeline.
Section~\ref{sec:lit} surveys related work.
Section~\ref{sec:data} describes the data.
Section~\ref{sec:ensemble} details the forecasting ensemble.
Section~\ref{sec:error} formalises the error characterisation.
Section~\ref{sec:confidence} derives the confidence scores and
Bayesian posterior.
Section~\ref{sec:sharpe} derives the Sharpe-optimal allocation
from first principles.
Section~\ref{sec:ot} derives the optimal transport execution.
Section~\ref{sec:results} presents the empirical results.

%======================================================================
\section{Related Work}\label{sec:lit}
%======================================================================

\textbf{Mean--variance optimisation and its fragility.}
Markowitz~\cite{markowitz1952} introduced mean--variance
optimisation; Sharpe~\cite{sharpe1966} proposed the reward-to-variability
ratio that bears his name.  The sensitivity of
$\bw^*=\Sigma^{-1}\bmu$ to input estimation error was analysed
by Merton~\cite{merton1980}, who showed that expected returns are
far harder to estimate than covariances from historical data.
Michaud~\cite{michaud1989} coined the term ``estimation error
maximiser'' for the Markowitz optimiser.  Best and
Grauer~\cite{best1991} proved that the tangent portfolio is
extremely sensitive to perturbations in the mean vector, with
the condition number of $\Sigma$ governing the amplification.
Chopra and Ziemba~\cite{chopra1993} quantified the relative
importance of errors in means versus variances versus covariances.
DeMiguel et~al.~\cite{demiguel2009} showed that the combined effect
of estimation error typically overwhelms the theoretical gains of
optimisation for sample sizes available in practice.

\textbf{Shrinkage and robust estimation.}
Ledoit and Wolf~\cite{ledoit2004} proposed shrinking the sample
covariance toward a structured target (constant correlation or
identity), choosing the intensity by minimising expected squared
Frobenius loss.  James and Stein~\cite{james1961} showed that
shrinking individual estimates toward a grand mean reduces total
expected squared error for $n\ge 3$.  Jorion~\cite{jorion1986}
applied this to portfolio means, showing that Bayes--Stein
shrinkage significantly improves out-of-sample performance.
Black and Litterman~\cite{black1992} embedded prior views through
a Bayesian framework, though their model requires the investor to
specify a view matrix and uncertainty.  Ben-Tal
et~al.~\cite{bental2009} formalised robust optimisation under
ellipsoidal uncertainty sets for both means and covariances.

\textbf{Machine learning for asset pricing.}
Gu et~al.~\cite{gu2020empirical} provided a comprehensive
comparison of ML methods (neural nets, trees, penalised linear
models) for cross-sectional stock return prediction, reporting
monthly $R^2$ up to 0.40\% out-of-sample.  Feng
et~al.~\cite{feng2020taming} used deep learning to discover
latent asset pricing factors.  Stacked
generalisation~\cite{wolpert1992,breiman1996} combines
heterogeneous learners through a meta-model, reducing variance at
modest bias cost; it is well-suited to noisy financial targets
because diversity among base learners prevents overfitting to any
single signal.

\textbf{Optimal transport in finance.}
Cuturi~\cite{cuturi2013} introduced entropic regularisation for
scalable OT computation; Peyr\'e and Cuturi~\cite{peyre2019}
provide a comprehensive survey.  In finance, Blanchet and
Murthy~\cite{blanchet2019} used Wasserstein ambiguity sets for
distributionally robust portfolio selection.  Mohajerin Esfahani
and Kuhn~\cite{esfahani2018} showed that Wasserstein DRO is
tractable and performs well for data-driven stochastic programming.
The application of OT to \emph{execution}---routing capital flows
from a current to a target portfolio at minimum cost---is, to our
knowledge, novel.

%======================================================================
\section{Data Pipeline}\label{sec:data}
%======================================================================

\subsection{Universe and horizon}

The asset universe consists of $n=24$ instruments: 19 global
equity indices spanning North America (S\&P~500, Nasdaq~100,
DJIA, Russell~2000, TSX), Europe (FTSE~100, DAX, Euro Stoxx~50,
CAC~40, IBEX~35), Asia-Pacific (Nikkei~225, Hang~Seng, ASX~200,
BSE Sensex, Nifty~50, KOSPI, TAIEX), and Latin America
(Bovespa, IPC Mexico); the VIX volatility index; and 4 commodity
futures (gold, silver, crude oil, natural gas).  Daily OHLCV data
is downloaded from Yahoo Finance (January 2003 -- December 2025)
with \texttt{auto\_adjust=False} to preserve raw OHLC ranges.
Training period: 2003--2011 (9 years).
Out-of-sample test: 2012--2025 (14 years).

\placeholder{Figure~1: Log-scale evolution of all 24 assets,
normalised to 1 at 2003 (Plot~P26).}

\subsection{Cleaning and resampling}

Per ticker: Close is forward-filled; Open/High/Low are
back-filled from Close if missing; Volume defaults to zero.
Tickers with fewer than 100 surviving daily bars are dropped;
$n$ is determined dynamically.  Daily data is aggregated to
weekly frequency anchored on Friday:
$O=\text{first}$, $H=\max$, $L=\min$, $C=\text{last}$,
$V=\sum$.  Weekly simple returns:
$r_{i,w}=C_{i,w}/C_{i,w-1}-1$.

\subsection{Volatility estimation}

Rather than close-to-close volatility (which uses only two
prices per period), we use the Garman--Klass
estimator~\cite{garman1980}, which exploits the full OHLC range:
\begin{equation}\label{eq:gk}
\hat\sigma_{i,w}^{\mathrm{GK}} = \sqrt{\max\!\Big(
  \tfrac{1}{2}(\ln H{-}\ln L)^2
  - (2\ln 2{-}1)(\ln C{-}\ln O)^2,\;0\Big)}.
\end{equation}
Under driftless geometric Brownian motion,
$\Var(\hat\sigma_{\mathrm{GK}}^2)\approx 0.273\,\sigma^4$,
yielding $\approx 7.3\times$ the statistical efficiency of
close-to-close estimation.  This matters because better
volatility estimates feed directly into a better covariance
matrix downstream.

%======================================================================
\section{Forecasting Ensemble}\label{sec:ensemble}
%======================================================================

Two prediction tasks are solved independently: next-week pairwise
correlation $\hat\rho_{ij,w+1}$ (one model per unordered pair) and
next-week asset return $r_{i,w+1}$ (one model per ticker).
Each task uses an identical stacked ensemble architecture.

\subsection{Feature construction}

\textbf{Correlation features} ($d_p=10$).
For each pair $(i,j)$ with $i<j$ and week $w$: four values of the
26-week rolling Pearson correlation at lags $\{0,1,2,4\}$, both
assets' current returns, both GK volatilities, their ratio
$\sigma_i/(\sigma_j+\epsilon)$, and their product $\sigma_i\sigma_j$.
The vol ratio and cross-vol features capture the empirical
correlation--volatility feedback documented by
Longin and Solnik~\cite{longin2001}: correlations rise during
high-volatility episodes, and asymmetrically so.

\textbf{Return features} ($d_r=14$).
For each ticker $i$ and week $w$: five lagged returns at horizons
$\{1,2,4,8,12\}$ weeks, three rolling means at 4-, 12-, and
26-week windows, two rolling standard deviations at 8 and 12 weeks,
the GK volatility, a vol-of-vol measure (standard deviation of GK
over 8 weeks), and a return acceleration term
$a=\bar r^{(4)}-\bar r^{(12)}$ that captures momentum changes.

\placeholder{Table~1: Feature summary with names, formulas,
and coverage statistics from the training set.}

\subsection{Stacked generalisation}

Stacked generalisation~\cite{wolpert1992}, also called stacked
regression~\cite{breiman1996}, combats overfitting by training a
meta-model on out-of-fold predictions of diverse base learners.
The theoretical motivation is bias--variance decomposition: if
base learners make uncorrelated errors, the meta-model's variance
decreases as $1/M$ while bias is controlled by the richest base
learner.

\textbf{Base learners} ($M=9$).
We use three model families for maximum diversity:
\begin{enumerate}
\item \textit{Linear models} (4): OLS (no regularisation), Ridge
  ($\alpha=1$, tunable in $[10^{-4},100]$), Lasso
  ($\alpha=10^{-3}$, \texttt{max\_iter}=5000), ElasticNet
  ($\alpha=10^{-3}$, $\ell_1$-ratio=0.5).
\item \textit{Tree ensembles} (4): Random Forest (250 trees,
  max depth 6, min leaf 5), Extra Trees (250 trees, depth 8,
  min leaf 3), gradient boosting (300 rounds, depth 3,
  $\eta=0.05$), AdaBoost (100 estimators, $\eta=0.05$).
\item \textit{Neural network} (1): a 3-hidden-layer MLP with 64
  units per layer, ReLU activations, dropout 0.1, trained with
  Adam ($\mathrm{lr}=10^{-3}$, weight decay $10^{-4}$) for up to
  100 epochs with early stopping (patience 12).
\end{enumerate}

Six of the nine (Ridge, Lasso, ElasticNet, RF, ExtraTrees, GBM)
are tuned via Optuna~\cite{akiba2019optuna} with TPE sampling,
10 trials or 20\,s wall-clock (whichever binds first), and 5-fold
\texttt{TimeSeriesSplit} cross-validation scored by MSE.

\textbf{Training procedure.}
For each pair or ticker: (i)~filter to the training period and
require $\ge 50$ observations; (ii)~split 80/20 chronologically
into base and meta sets; (iii)~fit a \texttt{StandardScaler} on
the base set and transform both; (iv)~train all 9 base learners
on the base set (with Optuna tuning where applicable);
(v)~predict on the meta set to form meta-features; (vi)~train the
meta-model; (vii)~store the scaler, base models, and meta-model.

\textbf{Meta-feature construction.}
After base training, predictions on the held-out meta set produce a
$2M$-dimensional vector per observation:
\begin{equation}\label{eq:meta}
\bm{z}_t = \big(\hat{y}_t^{(1)},\dots,\hat{y}_t^{(M)},\;
  |\hat{y}_t^{(1)}{-}y_t|,\dots,|\hat{y}_t^{(M)}{-}y_t|\big)
  \in\R^{18}.
\end{equation}
The first $M$ components are raw predictions; the second $M$ are
absolute errors.  A Ridge meta-model ($\alpha=1$) is trained on
$\bm{z}_t$ to learn the optimal combination.

At test time the absolute-error features are set to zero (the true
target is unavailable), so the meta-model reduces to a learned
linear blend of base predictions.  This is consistent with the code:
\texttt{np.hstack([bp, np.zeros\_like(bp)])}.

\textbf{Scale.}
The pipeline trains $\binom{n}{2}\approx 276$ pair ensembles and
$n\approx 24$ return ensembles.  Each ensemble has 9 base learners
and a meta-model, for a total of $\approx 2{,}700$ fitted models
per pair task and $\approx 240$ per return task.

\placeholder{Table~2: Per-base-learner MSE and $R^2$ on the meta
set, averaged across tickers, with ensemble (stacked) row.}

\placeholder{Figure~2: Meta-model Ridge coefficients for the return
task (Plot~P20), showing which base learners receive weight.}

%======================================================================
\section{Error Characterisation}\label{sec:error}
%======================================================================

Accurate predictions are necessary but not sufficient for good
portfolios.  What matters is knowing \emph{which} predictions to
trust.  At each test date $t$, the ensemble has been making
predictions for all prior weeks $s<t$.  We maintain a running
history of signed forecast errors
$\varepsilon_{i,s}=\hat r_{i,s}^{\mathrm{pred}}-r_{i,s+1}$
and distil five statistics per asset.

\begin{definition}[Return error matrix $E_\mu\in\R^{n\times 5}$]
\label{def:Emu}
Let $\mathcal{H}_{i,t}=\{\varepsilon_{i,s}\}_{s<t}$ and
$\lambda=0.94$ (half-life $\approx 11$ weeks).
\begin{align}
[E_\mu]_{i,1} &= \textstyle\sum_s w_s\varepsilon_{i,s}
  &&\text{(EWMA bias)},\label{eq:emu1}\\
[E_\mu]_{i,2} &= \mathbf{1}[\sign(\hat r_{i,t-1}){=}\sign(r_{i,t})]
  &&\text{(directional hit)},\label{eq:emu2}\\
[E_\mu]_{i,3} &= \mathrm{std}(\mathcal{H}_{i,t})
  &&\text{(error volatility)},\label{eq:emu3}\\
[E_\mu]_{i,4} &= \mathrm{skew}(\mathcal{H}_{i,t})
  &&\text{(error skewness)},\label{eq:emu4}\\
[E_\mu]_{i,5} &= \mathrm{exkurt}(\mathcal{H}_{i,t})
  &&\text{(error kurtosis)}.\label{eq:emu5}
\end{align}
\end{definition}

Column~1 detects \emph{systematic} over- or under-prediction, with
recent errors exponentially up-weighted.  Column~2 is a binary check:
even if the magnitude is wrong, does the model at least get the
\emph{sign} right?  Columns~3--5 characterise the \emph{shape} of
the error distribution.  A model whose errors are wide (high~3),
left-skewed (negative~4), and heavy-tailed (high~5) is
systematically dangerous: it over-predicts returns with occasional
large misses that the standard deviation underestimates.

An analogous tensor $E_C\in\R^{n\times n\times 4}$ tracks
correlation forecast errors per pair: EWMA bias, error std,
skewness, and kurtosis.

The last 26 weeks of GK volatility per asset are stored as
$V\in\R^{n\times 26}$, providing the raw material for
forward-looking covariance estimation.

\placeholder{Figure~3: Time series of EWMA bias and error vol for
5 representative tickers across the test period, showing
regime-dependent variation (source: error histories).}

%======================================================================
\section{Confidence and Bayesian Posterior}\label{sec:confidence}
%======================================================================

\subsection{From error statistics to confidence}

The five columns of $E_\mu$ answer a single question: \emph{how much
should we trust the return forecast for asset $i$?}  We compress
this into a scalar $\kappa_i\in(0,1]$ via four multiplicative
penalty terms, each designed to penalise a specific failure mode.

\textbf{Bias penalty.}
Large systematic error should kill trust.  We normalise by the
median absolute deviation (robust to outliers) and apply exponential
decay:
\begin{equation}\label{eq:sbias}
s_i^{\mathrm{bias}} = \exp\!\big({-}|b_i|/\MAD(\bm{b})\big),
\quad b_i=[E_\mu]_{i,1}.
\end{equation}

\textbf{Directional accuracy.}
A linear penalty: correct sign gives $s^{\mathrm{dir}}=1$;
a miss gives $0.7$.  Intentionally mild---one wrong sign
shouldn't destroy confidence:
\begin{equation}\label{eq:sdir}
s_i^{\mathrm{dir}} = 0.3\,d_i + 0.7,
\quad d_i=[E_\mu]_{i,2}\in\{0,1\}.
\end{equation}

\textbf{Dispersion penalty.}
High error volatility $v_i=[E_\mu]_{i,3}$ is dangerous, and
\emph{more} dangerous when the error distribution has fat tails
(excess kurtosis $k_i=[E_\mu]_{i,5}$), because then $v_i$
underestimates the true forecast risk:
\begin{equation}\label{eq:sdisp}
s_i^{\mathrm{disp}} = \exp\!\big({-}\delta_i/\mathrm{med}(\bm\delta)\big),
\quad \delta_i = v_i(1{+}\tfrac{1}{2}\max(k_i,0)).
\end{equation}

\textbf{Skewness penalty.}
Only \emph{negative} skew (over-prediction) is penalised---it is
asymmetrically harmful for a long-biased portfolio:
\begin{equation}\label{eq:sskew}
s_i^{\mathrm{skew}} = \exp\!\big({-}\max(-\xi_i,0)\big),
\quad \xi_i=[E_\mu]_{i,4}.
\end{equation}

\textbf{Composite confidence.}
The four scores combine as a geometric mean:
\begin{equation}\label{eq:kappa}
\kappa_i = \clip\!\Big(
  \prod_{k=1}^{4}(s_i^{(k)})^{1/4},\;10^{-4},\;1\Big).
\end{equation}
The multiplicative structure is deliberate.  An arithmetic mean
would allow a good directional score to mask a terrible bias;
the geometric mean ensures that \emph{any single} bad diagnostic
substantially reduces $\kappa_i$.

\textbf{Confidence-weighted returns.}
Each forecast is shrunk toward the cross-sectional mean
$\mu_0=n^{-1}\sum_j\hat\mu_j$:
\begin{equation}\label{eq:mu_tilde}
\tilde\mu_i = \kappa_i\hat\mu_i + (1{-}\kappa_i)\mu_0.
\end{equation}
When $\kappa_i=1$ we fully trust the forecast; when
$\kappa_i\to 0$ the asset's expected return converges to the
market average---maximum ignorance.  This is a form of
James--Stein shrinkage~\cite{james1961}, which provably reduces
expected squared error when $n\ge 3$.

\placeholder{Figure~4: $\kappa$ evolution for 10 tickers over the
test period (Plot~P13).}

\placeholder{Figure~5: Average $\kappa$ and average predicted vol,
dual axis (Plot~P14).}

\subsection{Bayesian posterior}

The confidence-weighted returns $\tilde\bmu$ are still point
estimates.  We embed them in a Normal--Inverse-Wishart (NIW)
conjugate model to produce a regularised mean and covariance.

\begin{definition}[NIW prior]
$\nu_0=n+2$ (the minimum for the IW mean to exist---a diffuse
prior), $\tau_0=1$ (prior weight equals one observation),
$\bm{m}_0=\mu_0\bm{1}$ (equal expected returns), and
$\Psi_0=\max(\nu_0-n-1,1)\hat\Sigma_{\mathrm{hist}}$, where
$\hat\Sigma_{\mathrm{hist}}$ is the historical covariance from
the training period.
\end{definition}

Treating $\tilde\bmu$ as a single observation modulated by the
confidence matrix $K=\diag(\bm\kappa)$, the posterior point
estimates are:
\begin{align}
\bmu^B &= \frac{\tau_0\bm{m}_0 + K\tilde\bmu}{\tau_0+1},
  \label{eq:muB}\\
\Sigma^B &= \frac{\Psi_0 + \frac{\tau_0}{\tau_0+1}
  (\tilde\bmu{-}\bm{m}_0)(\tilde\bmu{-}\bm{m}_0)^\top}
  {\max(\nu_0{+}1{-}n{-}1,1)}.
  \label{eq:SigB}
\end{align}
Both are projected to the nearest positive-definite matrix via
eigenvalue clamping~\cite{higham2002}.

The posterior mean $\bmu^B$ interpolates between the prior
$\mu_0\bm{1}$ and the data $\tilde\bmu$, weighted by
$\tau_0$ versus confidence.  When all $\kappa_i$ are small,
$\bmu^B\approx\mu_0\bm{1}$---the model admits it knows nothing
special.  When all $\kappa_i\approx 1$ and $\tau_0$ is small,
$\bmu^B\approx\hat\bmu$---the data speaks.

\subsection{Final covariance}

EWMA volatility from $V$ with decay $\lambda_v=0.94$ and
half-life $\approx 11$ weeks:
$\hat\sigma_i = \sqrt{\sum_\ell w_\ell V_{i,\ell}^2}$.
A regime adjustment inflates volatility when current levels
exceed their own 26-week average:
$\hat\sigma_i^{\mathrm{adj}}=\hat\sigma_i(1+0.5\max(r_i^{\mathrm{vol}}-1,0))$,
where $r_i^{\mathrm{vol}}=\hat\sigma_i/\bar V_i$.  This is a
one-sided adjustment: vol is only inflated, never deflated,
biasing toward conservatism during stress.

The final covariance blends Bayesian and forward-looking estimates:
\begin{equation}\label{eq:Sigma_final}
\Sigma_{\mathrm{final}} = \tfrac{1}{2}\Sigma^B
  + \tfrac{1}{2} D_\sigma\hat C_{t+1}D_\sigma,
\end{equation}
where $D_\sigma=\diag(\hat\sigma_1^{\mathrm{adj}},\dots,
\hat\sigma_n^{\mathrm{adj}})$ and $\hat C_{t+1}$ is the
predicted correlation matrix.  The 50/50 blend avoids
putting full weight on either the backward-looking Bayesian
estimate or the forward-looking ML prediction.

\placeholder{Figure~6: Heatmap of $\bmu^B$ (Plot~P15) and
$\hat\sigma^{\mathrm{adj}}$ (Plot~P16) over the test period.}

%======================================================================
\section{Sharpe-Optimal Allocation}\label{sec:sharpe}
%======================================================================

\subsection{Why Sharpe maximisation?}

Among the many portfolio objectives (minimum variance, maximum
diversification, risk parity), Sharpe maximisation targets the
tangent portfolio on the efficient frontier---the portfolio offering
the highest expected excess return per unit of risk.  Under the
assumption of a single risk-free rate, the tangent portfolio is the
unique intersection of the capital market line with the frontier,
and any investor's optimal allocation is a combination of the
risk-free asset and the tangent portfolio (Tobin's two-fund
separation theorem).

\subsection{From fractional programme to QP sweep}

The Sharpe ratio is:
\begin{equation}\label{eq:SR}
\mathrm{SR}(\bw) = \frac{\bw^\top\bmu^B - r_f}
  {\sqrt{\bw^\top\Sigma_{\mathrm{final}}\bw}},
\end{equation}
where $r_f=r_f^{\mathrm{ann}}/52$ is the weekly risk-free rate.
Maximising~\eqref{eq:SR} directly is a quasi-convex fractional
programme.  The standard approach~\cite{cornuejols2007} is to note
that for any $\gamma>0$, maximising the risk-adjusted return
\begin{equation}\label{eq:J}
J(\bw;\gamma) = (\bmu^B{-}r_f\bm{1})^\top\bw
  - \gamma\,\bw^\top\Sigma_{\mathrm{final}}\bw
\end{equation}
subject to a constraint set yields a portfolio on the efficient
frontier.  As $\gamma$ varies from $0$ to $\infty$, the solution
traces the frontier from the maximum-return to the minimum-variance
portfolio.  The $\gamma$ that achieves the highest Sharpe identifies
the tangent portfolio.

\subsection{KKT derivation}

To build intuition, consider first the case with only a budget
constraint $\bm{1}^\top\bw=1$.  The Lagrangian is:
\begin{equation}\label{eq:lagrangian}
\mathcal{L} = (\bmu^B{-}r_f\bm{1})^\top\bw
  - \gamma\bw^\top\Sigma_{\mathrm{final}}\bw
  - \lambda(\bm{1}^\top\bw-1).
\end{equation}
Setting $\nabla_{\bw}\mathcal{L}=0$:
\begin{equation}\label{eq:foc}
\bmu^B - r_f\bm{1} - 2\gamma\Sigma_{\mathrm{final}}\bw
  - \lambda\bm{1} = 0.
\end{equation}
Solving for $\bw$:
\begin{equation}\label{eq:w_closed}
\bw = \frac{1}{2\gamma}\Sigma_{\mathrm{final}}^{-1}
  \big(\bmu^B - r_f\bm{1} - \lambda\bm{1}\big).
\end{equation}
Substituting into $\bm{1}^\top\bw=1$ determines $\lambda$.
Write $A=\bm{1}^\top\Sigma_{\mathrm{final}}^{-1}\bm{1}$ and
$B=\bm{1}^\top\Sigma_{\mathrm{final}}^{-1}(\bmu^B-r_f\bm{1})$:
\begin{equation}\label{eq:lambda}
\lambda = \frac{B - 2\gamma}{A}.
\end{equation}
Substituting back:
\begin{equation}\label{eq:w_final}
\bw^*(\gamma) = \frac{1}{2\gamma}\Sigma_{\mathrm{final}}^{-1}
  \Big(\bmu^B - r_f\bm{1}
  - \frac{B-2\gamma}{A}\bm{1}\Big).
\end{equation}
This is the classical two-fund separation result.  The optimal
portfolio is a linear combination of
$\Sigma_{\mathrm{final}}^{-1}(\bmu^B-r_f\bm{1})$
(the ``tangent direction'') and
$\Sigma_{\mathrm{final}}^{-1}\bm{1}$
(the minimum-variance direction), with mixing controlled by $\gamma$.

\begin{remark}[Why the closed form is dangerous]
Equation~\eqref{eq:w_final} inverts $\Sigma_{\mathrm{final}}$ and
multiplies by $\bmu^B$.  If either contains noise, $\Sigma^{-1}$
amplifies it---this is precisely the Michaud
critique~\cite{michaud1989}.  Our remedy operates upstream:
$\bmu^B$ has been shrunk by confidence-driven Bayesian updating
(Section~\ref{sec:confidence}), and $\Sigma_{\mathrm{final}}$ is
a regularised blend (Eq.~\ref{eq:Sigma_final}).  The optimiser
sees de-noised inputs whose estimation error has been
systematically attenuated.
\end{remark}

\subsection{Constrained programme}

In practice we add leverage and concentration constraints that
further regularise the solution.  The full programme for each
$\gamma$ is:
\begin{equation}\label{eq:qp_full}
\bw^*(\gamma) = \arg\max_{\bw\in\R^n}\;J(\bw;\gamma)
\end{equation}
subject to:
\begin{align}
\bm{1}^\top\bw &= 1
  &&\text{(fully invested)},\label{eq:c_budget}\\
\|\bw\|_1 &\le 1.6
  &&\text{(gross leverage)},\label{eq:c_lev}\\
|w_i| &\le 0.20,\;\forall\,i
  &&\text{(concentration)}.\label{eq:c_conc}
\end{align}
In long-only mode, $w_i\ge 0$ replaces the lower bound.  Each
sub-problem is a convex QP ($\Sigma_{\mathrm{final}}\succ 0$)
solved via SCS (\texttt{max\_iters}=5000).

We sweep $\gamma$ over 50 log-spaced values in $[10^{-2},10^3]$
and select:
\begin{equation}\label{eq:gamma_star}
\gamma^* = \arg\max_{\gamma\in\Gamma}\;\mathrm{SR}(\bw^*(\gamma)),
\qquad \bw^*_{t+1} = \bw^*(\gamma^*).
\end{equation}

The leverage constraint $\|\bw\|_1\le 1.6$ permits moderate
shorting (net leverage up to 1.6$\times$), while the box constraint
$|w_i|\le 0.20$ prevents concentration in a single asset.  Together
these act as implicit regularisers that prevent the optimiser from
taking extreme positions even when one asset's $\mu^B_i$ is an
outlier.

\placeholder{Figure~7: Weight stacks (Plot~P07) and weight
heatmap (Plot~P22) over the test period.}

\placeholder{Figure~8: Gross and net leverage (Plot~P08).}

\subsection{Volatility targeting}

Post-allocation, portfolio returns are scaled to a target annualised
volatility $\sigma^*=0.15$:
\begin{equation}\label{eq:vt}
r_t^{\mathrm{VT}} = r_t^{\mathrm{raw}}\cdot
  \min\!\Big(\frac{\sigma^*}
  {\max(\hat\sigma_t^{\mathrm{port}},10^{-6})},\;2.5\Big),
\end{equation}
where $\hat\sigma_t^{\mathrm{port}}=
\mathrm{std}(\{r_s^{\mathrm{raw}}\}_{s=t-25}^{t})\sqrt{52}$.
During calm periods leverage is increased (up to $2.5\times$);
during crises exposure is cut.  All strategies (PRISM,
equal-weight, SPY) are vol-fitted identically, enabling fair
comparison at the same risk budget.

%======================================================================
\section{Optimal Transport Execution}\label{sec:ot}
%======================================================================

Given the current portfolio $\bw_t$ and target $\bw^*_{t+1}$, we
must rebalance.  The na\"ive approach computes net trades
$\Delta w_i=w_i^*-w_{t,i}$ and executes them independently.
Optimal transport asks: \emph{what is the cheapest way to rearrange
the current weights into the target, accounting for the pairwise
structure of transaction costs?}

\subsection{Kantorovich formulation}

Normalise to probability distributions:
$\bm{a}=|\bw_t|/\||\bw_t|\|_1$,
$\bm{b}=|\bw^*|/\||\bw^*|\|_1$.  The cost matrix has entries
$C_{ij}=|c_i|+|c_j|$ for $i\neq j$ (10\,bps per leg),
$C_{ii}=0$.  The transport polytope is
$\Pi(\bm{a},\bm{b})=\{T\ge 0:T\bm{1}=\bm{a},\,T^\top\bm{1}=\bm{b}\}$.
The Kantorovich problem is:
\begin{equation}\label{eq:kant}
T^* = \arg\min_{T\in\Pi(\bm{a},\bm{b})}
  \langle C,T\rangle.
\end{equation}

\subsection{Entropic regularisation and Sinkhorn}

Adding $\varepsilon\sum_{ij}T_{ij}(\log T_{ij}-1)$ to the
objective, with $\varepsilon=0.05\cdot\mathrm{median}(C)$, makes
the problem strictly convex.

\begin{proposition}[Solution structure]
The first-order conditions give
$C_{ij}+\varepsilon\log T_{ij}-f_i-g_j=0$, hence
$T_{ij}^*=u_i K_{ij}v_j$ where $K_{ij}=e^{-C_{ij}/\varepsilon}$,
$u_i=e^{f_i/\varepsilon}$, $v_j=e^{g_j/\varepsilon}$.
The scaling vectors are determined by the marginal constraints.
\end{proposition}

The Sinkhorn algorithm finds $\bm{u},\bm{v}$ by alternating
projections: $\bm{u}\leftarrow\bm{a}\oslash(K\bm{v})$,
$\bm{v}\leftarrow\bm{b}\oslash(K^\top\bm{u})$.  Each iteration
is a coordinate ascent step on the strictly concave dual; convergence
is linear at rate $O(e^{-c/\varepsilon})$, typically 10--30
iterations at tolerance $10^{-8}$.

\textbf{Transaction cost.}
$\mathcal{C}^*=\sum_{i\neq j}C_{ij}T^*_{ij}$.
Net return: $r_t=(\bw^*)^\top\bm{r}_t-\mathcal{C}^*$.

\textbf{Routing.}
The net trades $w_i^*-w_{t,i}$ are identical to na\"ive
rebalancing by the marginal constraints.  The value of OT is in
the pairwise routing: $T^*_{ij}$ tells the execution desk which
sells fund which buys, enabling pair or basket execution.

A synthetic order book (mid = Close, spread $\approx$ GK vol
$\times 0.1$, five depth levels) logs every flow $>10^{-8}$
with estimated slippage.

\placeholder{Figure~9: Net trades from the last rebalance
(Plot~P31).  Turnover and TC bars (Plot~P09).}

\placeholder{Table~3: Execution summary---avg turnover, avg TC,
total fills, total slippage.}

%======================================================================
\section{Results}\label{sec:results}
%======================================================================

\subsection{Forecast quality}

\placeholder{Table~4: Per-ticker return model $R^2$, MAE, DA.
All tickers sorted by $R^2$.}

\placeholder{Table~5: Aggregate pair model diagnostics---median
$R^2$, DA, MAE across all pairs.}

\placeholder{Figure~10: Return model $R^2$ and DA bar charts
(Plots~P18, P19).}

Key questions: Is $R^2$ positive for the majority?  Is DA
$>0.5$?  Does stacking beat the best base learner?

\subsection{Portfolio performance}

\placeholder{Table~6: Full metrics for PRISM Raw, PRISM VT,
PRISM @15\%, EW @15\%, SPY @15\%.  Columns: Total, Ann.\ Return,
Ann.\ Vol, Sharpe, Sortino, MDD, Calmar, Win\%, Skew, Kurt.}

\placeholder{Figure~11: Vol-matched cumulative returns at 15\%,
log scale (Plot~P02).}

\placeholder{Figure~12: Drawdowns (Plot~P03).  Rolling Sharpe
(Plot~P04).  Monthly heatmap (Plot~P12).  Annual returns
(Plot~P23).}

\subsection{Confidence diagnostics}

\placeholder{Figure~13: Avg $\kappa$ vs predicted vol (Plot~P14).
Hypothesis: inverse relationship; $\kappa$ drops during GFC,
COVID, 2022.}

\placeholder{Figure~14: Rolling vol with 15\% target (Plot~P06).}

\subsection{Ablation}

\placeholder{Table~7: Sharpe, MDD, Calmar for: (a) raw Markowitz,
(b) +confidence, (c) full pipeline.}

%======================================================================
\section{Discussion}\label{sec:discussion}
%======================================================================

\textbf{Error characterisation bridges forecasting and allocation.}
The five-column $E_\mu$ forces the allocator to respect forecast
uncertainty.  Assets with high bias, low directional accuracy, or
fat-tailed errors are de-weighted automatically, without
hand-tuned position limits.

\textbf{The Sharpe derivation justifies the $\gamma$-sweep.}
The KKT analysis shows that the unconstrained solution
$\bw\propto\Sigma^{-1}\bmu$ is dangerously sensitive; the
regularised inputs and convex constraints jointly tame this.
The $\gamma$-sweep recovers the tangent portfolio on the
regularised frontier.

\textbf{Comparison with Ledoit--Wolf and Black--Litterman.}
Ledoit--Wolf~\cite{ledoit2004} shrinks the covariance with a fixed
intensity; PRISM adapts shrinkage per-asset via $\kappa_i$.
Black--Litterman~\cite{black1992} requires the investor to specify
views; PRISM derives its ``views'' from the ML forecasts and
its ``view uncertainty'' from the error characterisation.

\textbf{OT provides structural insight.}
Beyond cost, the transport plan reveals which assets cross-fund each
other---information invisible in net trades.

\textbf{Limitations.}
(1)~Static training (no walk-forward retraining).
(2)~Synthetic order book.
(3)~Meta-model error features zeroed at test time.
(4)~Survivorship bias from a fixed universe.
(5)~The blend weight $\zeta=0.5$ is not optimised.

\textbf{Extensions.}
Walk-forward retraining; macro features; conformal prediction
intervals replacing the heuristic $\kappa$; GNN base learner on
the correlation graph; multi-asset-class universes.

%======================================================================
\section{Conclusion}\label{sec:conclusion}
%======================================================================

PRISM demonstrates that the gap between ML forecasts and portfolio
construction can be closed by treating prediction uncertainty as a
first-class object.  The five-tensor interface between forecasting
and allocation forces the optimiser to acknowledge what the model
does and does not know.  Confidence-driven Bayesian shrinkage
de-noises expected returns; robust covariance blending regularises
the risk model; the KKT-derived $\gamma$-sweep identifies the
tangent portfolio on the regularised frontier; and Sinkhorn optimal
transport executes the result at minimum cost.  Every equation maps
to the implemented codebase.

%======================================================================
\newpage
{\small
\begin{thebibliography}{99}

\bibitem{markowitz1952}
H.~Markowitz, ``Portfolio selection,''
\textit{J.~Finance}, 7(1):77--91, 1952.

\bibitem{sharpe1966}
W.F.~Sharpe, ``Mutual fund performance,''
\textit{J.~Business}, 39(1):119--138, 1966.

\bibitem{merton1980}
R.C.~Merton, ``On estimating the expected return on the market,''
\textit{J.~Financial Econ.}, 8(4):323--361, 1980.

\bibitem{michaud1989}
R.O.~Michaud, ``The Markowitz optimization enigma,''
\textit{Financial Analysts J.}, 45(1):31--42, 1989.

\bibitem{best1991}
M.J.~Best, R.R.~Grauer, ``On the sensitivity of mean-variance-efficient
portfolios to changes in asset means,''
\textit{Rev.~Financial Studies}, 4(2):315--342, 1991.

\bibitem{chopra1993}
V.K.~Chopra, W.T.~Ziemba, ``The effect of errors in means, variances,
and covariances on optimal portfolio choice,''
\textit{J.~Portfolio Management}, 19(2):6--11, 1993.

\bibitem{demiguel2009}
V.~DeMiguel, L.~Garlappi, R.~Uppal, ``Optimal versus naive
diversification,''
\textit{Rev.~Financial Studies}, 22(5):1915--1953, 2009.

\bibitem{james1961}
W.~James, C.~Stein, ``Estimation with quadratic loss,''
\textit{Proc.~4th Berkeley Symp.}, 1:361--379, 1961.

\bibitem{jorion1986}
P.~Jorion, ``Bayes-Stein estimation for portfolio analysis,''
\textit{J.~Financial Quant.~Anal.}, 21(3):279--292, 1986.

\bibitem{ledoit2004}
O.~Ledoit, M.~Wolf, ``A well-conditioned estimator for
large-dimensional covariance matrices,''
\textit{J.~Multivariate Anal.}, 88(2):365--411, 2004.

\bibitem{black1992}
F.~Black, R.~Litterman, ``Global portfolio optimization,''
\textit{Financial Analysts J.}, 48(5):28--43, 1992.

\bibitem{bental2009}
A.~Ben-Tal, L.~El~Ghaoui, A.~Nemirovski,
\textit{Robust Optimization}, Princeton Univ.~Press, 2009.

\bibitem{gu2020empirical}
S.~Gu, B.~Kelly, D.~Xiu, ``Empirical asset pricing via machine
learning,''
\textit{Rev.~Financial Studies}, 33(5):2223--2273, 2020.

\bibitem{feng2020taming}
G.~Feng, S.~Giglio, D.~Xiu, ``Taming the factor zoo,''
\textit{J.~Finance}, 75(3):1327--1370, 2020.

\bibitem{longin2001}
F.~Longin, B.~Solnik, ``Extreme correlation of international equity
markets,''
\textit{J.~Finance}, 56(2):649--676, 2001.

\bibitem{wolpert1992}
D.H.~Wolpert, ``Stacked generalization,''
\textit{Neural Networks}, 5(2):241--259, 1992.

\bibitem{breiman1996}
L.~Breiman, ``Stacked regressions,''
\textit{Machine Learning}, 24(1):49--64, 1996.

\bibitem{garman1980}
M.B.~Garman, M.J.~Klass, ``On the estimation of security price
volatilities from historical data,''
\textit{J.~Business}, 53(1):67--78, 1980.

\bibitem{cuturi2013}
M.~Cuturi, ``Sinkhorn distances,''
\textit{Advances in NIPS}, 26, 2013.

\bibitem{peyre2019}
G.~Peyr\'e, M.~Cuturi, ``Computational Optimal Transport,''
\textit{Found.~Trends ML}, 11(5--6):355--607, 2019.

\bibitem{blanchet2019}
J.~Blanchet, K.~Murthy, ``Quantifying distributional model risk
via optimal transport,''
\textit{Math.~Oper.~Res.}, 44(2):565--600, 2019.

\bibitem{esfahani2018}
P.~Mohajerin~Esfahani, D.~Kuhn, ``Data-driven distributionally
robust optimization using the Wasserstein metric,''
\textit{Math.~Programming}, 171:115--166, 2018.

\bibitem{cornuejols2007}
G.~Cornuejols, R.~T\"ut\"unc\"u,
\textit{Optimization Methods in Finance}, Cambridge Univ.~Press, 2007.

\bibitem{higham2002}
N.J.~Higham, ``Computing the nearest correlation matrix,''
\textit{IMA J.~Numer.~Anal.}, 22(3):329--343, 2002.

\bibitem{akiba2019optuna}
T.~Akiba et al., ``Optuna: A next-generation hyperparameter
optimization framework,''
\textit{Proc.~25th ACM SIGKDD}, 2623--2631, 2019.

\end{thebibliography}
}

\end{document}